\documentclass[12pt]{article}
\usepackage{ctex}
\usepackage{amsmath,amssymb}
\usepackage{graphicx}
\usepackage{hyperref}   %生成超链接
\usepackage{enumitem}
\usepackage{amsfonts}   %数学字体
\usepackage{array}
\usepackage{booktabs}
\usepackage{xltabular}    %跨页+自适应宽度=longtable+tabularx

\hypersetup{
    colorlinks=true,   %去掉链接边框，通过文字颜色识别链接
    linkcolor=black,   %内部链接为黑色
    urlcolor=black,    %网址文字为黑色
    citecolor=black,   %参考文献引用标注为黑色
}

\begin{document}
\title{机器学习总结汇报}
\author{lyy}
\date{2026年7月30日}
\maketitle
%摘要
\begin{abstract}   
本文从理论与数学层面梳理机器学习入门核心知识点，一方面介绍机器学习的基本原理，另一方面对深度学习网络架构、传统机器学习模型进行分类解析。
\end{abstract}
%目录
\tableofcontents
\newpage


\section{机器学习基本原理}

\subsection{机器学习(Machine Learning)}
\subsubsection{机器学习定义}
机器学习是一种通过数据训练模型，使模型能够从数据中学习规律并做出预测或决策的技术。它是人工智能的核心领域之一，涵盖了多种算法和方法，如监督学习、无监督学习和半监督学习。
\subsubsection{数学表达}
传统机器学习基于静态标注数据集学习输入输出映射关系，依赖人工特征工程完成特征构建。
给定数据集 $\mathcal{D}=\{(x_i,y_i)\}_{i=1}^N$，模型参数为 $\boldsymbol{w}$，预测输出为 $\hat{y}=f(x;\boldsymbol{w})$。
机器学习核心优化目标为最小化经验损失：
\[
\boldsymbol{w}^*=\mathop{\arg\min}_{\boldsymbol{w}} \frac{1}{N}\sum_{i=1}^N L\big(f(x_i;\boldsymbol{w}),y_i\big)
\]

\subsubsection{应用场景}
\begin{itemize}
    \item 图像识别：在人脸识别系统中，通过收集大量不同人的面部图像数据，利用机器学习算法训练模型，使模型学习到不同人脸的特征模式。当新的人脸图像输入时，模型可以快速准确地识别出该人脸属于哪个人，广泛应用于安防监控、门禁系统等领域。
    
    \item 自然语言处理：在情感分析方面，通过分析用户在社交媒体、电商平台等留下的评论数据，判断用户对产品或服务的情感倾向是积极、消极还是中性，有助于企业了解用户反馈，改进产品和服务。
    
    \item 推荐系统：电商平台根据用户的历史购买记录、浏览行为、搜索关键词等数据，运用机器学习算法分析用户的兴趣偏好，为用户推荐可能感兴趣的商品，提高用户的购买转化率和平台的销售额，如淘宝、京东等电商平台的个性化推荐功能。
    
    \item 数据挖掘：在客户分群中，通过分析客户的年龄、性别、消费习惯、购买频率等多维度数据，使用聚类算法将客户分成不同的群体，企业可以针对不同群体制定个性化的营销策略，提高营销效果。

\end{itemize}


\subsection{深度学习(Deep Learning)}
\subsubsection{深度学习定义}
深度学习是机器学习的一个子领域，使用多层神经网络（深度神经网络）来学习数据的复杂特征。深度学习通过模拟人脑的神经网络结构，能够处理高维和非线性数据。

\subsubsection{数学表达}
深度学习是机器学习的子集，依托多层神经网络实现自动特征提取，通过多层非线性复合函数拟合高维复杂数据分布。
深层网络数学表达为多层嵌套函数：
\[
f(x;\boldsymbol{w})=f_L\big(f_{L-1}(\dots f_1(x;\boldsymbol{w}_1);\boldsymbol{w}_2\dots);\boldsymbol{w}_L\big)
\]
单层神经网络前向运算公式：
\[
\boldsymbol{h}^{(l)}=\sigma\big(\boldsymbol{W}^{(l)}\boldsymbol{h}^{(l-1)}+\boldsymbol{b}^{(l)}\big)
\]
其中 $\sigma(\cdot)$ 为激活函数，引入非线性能力，是深层模型拟合复杂任务的核心。

\subsubsection{应用场景}
\begin{itemize}
    \item 计算机视觉：在目标检测方面，像 Faster R-CNN、YOLO 系列算法可以快速检测出图像或视频中的目标物体，并确定其位置和类别，常用于智能交通中的车辆检测和计数，以及自动驾驶中的行人、车辆、交通标志等检测。

    \item 自然语言处理：机器翻译借助 Transformer 架构的神经网络，如谷歌的神经机器翻译系统，能够实现不同语言之间的自动翻译，打破语言障碍，方便跨国交流和信息传播。

    \item 语音处理：语音助手利用深度学习技术实现语音识别、语义理解和语音合成等功能，如小爱同学、天猫精灵等，能够理解用户的语音指令并提供相应的服务，如查询信息、播放音乐、控制智能家居设备等。

    \item 自动驾驶：深度学习在自动驾驶中起着关键作用，通过对摄像头采集的道路图像进行分析，识别道路边界、车道线、交通标志和信号灯等。

\end{itemize}

\subsection{强化学习(Reinforcement Learning)}
\subsubsection{强化学习定义}
强化学习是一种通过与环境交互学习策略的技术。智能体通过试错和奖励机制来优化行为策略，目标是最大化累积奖励。

\subsubsection{数学表达}
强化学习不同于监督学习，基于马尔可夫决策过程，通过智能体与环境交互试错学习最优策略，无需标注数据。
定义折扣累积回报：
\[
G_t=\sum_{k=0}^{\infty}\gamma^k r_{t+k},\quad 0<\gamma\le1
\]
优化目标为最大化长期期望回报：
\[
\pi^*=\mathop{\arg\max}_{\pi}\mathbb{E}[G_t]
\]

\subsubsection{应用场景}
\begin{itemize}
    \item 游戏 AI：强化学习在游戏领域取得了举世瞩目的成果，最为著名的当属 AlphaGo。AlphaGo 结合了深度学习和强化学习技术，通过自我对弈进行训练，学习到了高超的围棋策略，成功击败了人类围棋世界冠军，震惊了全世界。这一突破展示了强化学习在解决复杂策略性问题上的巨大潜力，也为游戏 AI 的发展开辟了新的道路。

    \item 机器人控制：在机器人控制领域，强化学习可以使机器人根据环境的变化自主学习最优的控制策略。在机械臂控制中，强化学习算法可以根据机械臂的当前位置、目标位置以及周围环境信息，学习如何精确地控制机械臂的关节运动，完成抓取、装配等复杂任务。

    \item 自动驾驶：强化学习在自动驾驶中发挥着关键作用，可用于车辆的路径规划、速度控制和动态决策等任务。通过传感器获取周围环境的信息，如道路状况、交通信号、其他车辆和行人的位置等，自动驾驶系统利用强化学习算法学习如何根据这些信息做出最优的驾驶决策，以实现安全、高效的行驶。

    \item 资源调度：在网络资源分配中，强化学习可以根据网络流量、用户需求等动态信息，学习如何合理分配带宽、服务器资源等，以提高网络性能和用户体验。

\end{itemize}

\subsection{三者对比}
\begin{xltabular}{\textwidth}{XXXX}
\caption{机器学习、深度学习与强化学习对比} \\
\toprule  %三线表的顶部粗线
\textbf{特性} & \textbf{机器学习 (ML)} & \textbf{深度学习 (DL)} & \textbf{强化学习 (RL)} \\
\midrule  %表头下方细线
\endfirsthead
\caption[]{机器学习、深度学习与强化学习对比（续）} \\
\toprule
\textbf{特性} & \textbf{机器学习 (ML)} & \textbf{深度学习 (DL)} & \textbf{强化学习 (RL)} \\
\midrule
\endhead
\bottomrule  %画出底部粗线
\endfoot
定义 & 从数据中学习规律并做出预测或决策 & 使用深度神经网络学习数据的复杂特征 & 通过与环境的交互学习策略 \\
数据需求 & 需要大量标注数据（监督学习） & 需要大量标注数据 & 不需要标注数据，但需要环境交互 \\
模型复杂度 & 模型相对简单（如线性回归、决策树） & 模型复杂（如深度神经网络） & 模型复杂度取决于任务和环境 \\
训练方式 & 基于数据训练 & 基于数据训练 & 基于环境交互和奖励机制 \\
应用场景 & 图像识别、推荐系统、数据挖掘 & 计算机视觉、自然语言处理、语音识别 & 游戏 AI、机器人控制、自动驾驶 \\
优点 & 易于理解和实现 & 能够处理复杂任务和高维数据 & 能够处理动态环境和序列决策问题 \\
缺点 & 对复杂任务表现有限 & 需要大量计算资源和数据 & 训练过程不稳定，收敛速度慢 \\
\end{xltabular}

\subsection{模型训练过程}
\subsubsection{准备工作}
模型进行训练前，需要进行数据预处理、构建模型、设置损失函数和优化器等工作。
\begin{itemize}
    \item 数据预处理：数据即为我们的研究对象，所有数据首先要经过数据清洗、归一化、缺失值处理等预处理，然后将数据集按照一定比例分为训练集、验证集和测试集。 
    \item 构建模型: 根据任务的需求和数据的特点，选择合适的模型结构，模型结构可以简化为我们学习过的\textbf{f(x)}函数，即带参数的数学函数，参数的数量和连接方式决定了模型的复杂程度和能力。    
    \item 设置损失函数和优化器：损失函数用于衡量模型预测结果与真实结果之间的差异，优化器则负责调整模型的参数，以最小化损失函数的值。
\end{itemize}

\subsubsection{核心步骤}
模型训练是一个反复迭代的过程，其固定的执行前向传播、计算损失、反向传播、参数更新四个步骤。
\begin{itemize}
    \item 前向传播：数据从输入端开始，依次通过模型的每层计算，最终得到预测结果，该过程只负责计算输出，不会更新参数。
    \item 计算损失：根据预测输出和真实标签计算损失值。
    \item 反向传播：将得到的损失从输出向输入逐层传播，计算每一层参数对损失的贡献（梯度），有助于模型调整参数，缩小损失，提高模型的性能。
    \item 参数更新：根据计算得到的梯度，使用优化算法更新模型的参数。
    
\end{itemize}

\subsubsection{收敛}
在优化领域，收敛通常指数值迭代序列趋于某一稳定点的过程。
针对机器学习训练，可从两个层面理解：
\begin{enumerate}
    \item \textbf{参数收敛}：参数更新幅度趋于零。
    \item \textbf{目标收敛}：训练损失或验证损失趋于平稳，不再显著下降。
\end{enumerate}
实际训练过程中，我们可以通过观察损失曲线是否逐渐趋于平坦来判断模型是否收敛。
\subsubsection{训练过程中的监控与控制}
模型在训练过程中需要持续观察和干预，不是每次都要等待模型训练结束才能评估模型训练的状态，下面介绍一些监控模型训练的方法。
\begin{itemize}
    \item 监测损失值：可以查看.csv文件，训练集和验证集的损失在正常情况下都应该逐渐下降。
    \item 过拟合判断：当训练集损失继续下降，但验证集损失反而上升时，说明模型对新数据的预测能力变差，发生了过拟合。
    \item 早停：在验证集损失不再下降甚至上升时，提前终止训练，避免浪费时间和过拟合，也可以设置早停的超参数，当损失不再下降，并超过一定轮数，模型会自动停止训练。
    \item 学习率调整：当模型学习过程振荡上升时，可以调整模型的学习率，适当降低，让模型学习的更精细，逐步接近最优。
    \item 正则化：在损失函数中加入惩罚项，限制模型不要过度依赖某些参数，学习更通用的特征，从而提升泛化能力。
\end{itemize}

\subsection{基础衡量指标}
评估模型性能的指标因任务类型而异。
\subsubsection{分类任务}
记真正例 (TP)、假正例 (FP)、真负例 (TN)、假负例 (FN)。常用指标如下：
\begin{itemize}
    \item \textbf{准确率}：$\displaystyle \text{Accuracy}=\frac{TP+TN}{TP+FP+TN+FN}$
    \item \textbf{精确率}：$\displaystyle \text{Precision}=\frac{TP}{TP+FP}$
    \item \textbf{召回率}：$\displaystyle \text{Recall}=\frac{TP}{TP+FN}$
    \item \textbf{F1 分数}：$\displaystyle \text{F1}=2\cdot\frac{\text{Precision}\cdot\text{Recall}}{\text{Precision}+\text{Recall}}$
\end{itemize}

\subsubsection{回归任务}
设真实值为 $y_i$，预测值为 $\hat{y}_i$，$N$ 为样本数。
\begin{itemize}
    \item \textbf{均方误差}：$\displaystyle \text{MSE}=\frac{1}{N}\sum_{i=1}^{N}(y_i-\hat{y}_i)^2$
    \item \textbf{平均绝对误差}：$\displaystyle \text{MAE}=\frac{1}{N}\sum_{i=1}^{N}|y_i-\hat{y}_i|$
    \item \textbf{决定系数}：$\displaystyle R^2=1-\frac{\sum_i(y_i-\hat{y}_i)^2}{\sum_i(y_i-\bar{y})^2}$，其中 $\bar{y}$ 为均值。
\end{itemize}





\section{基本模型解析}
%深度学习
\subsection{深度学习(Deep Learning)}
\subsubsection{神经网络(Neural Network)}
\paragraph{多层感知机(MLP)}
\subparagraph{模型概念}
多层感知机（Multilayer Perceptron, MLP）是最基础的前馈神经网络，由输入层、若干隐藏层和输出层全连接构成，每层神经元与下一层所有神经元两两相连，因此也被称为全连接神经网络。
\begin{itemize}
    \item \textbf{核心结构}：输入层接收原始特征数据；隐藏层通过线性变换+非线性激活函数对特征进行逐层抽象，常用激活函数包括ReLU、Sigmoid、Tanh等；输出层根据任务类型输出结果（回归任务输出连续值，分类任务通过Softmax输出类别概率）。
    \item \textbf{训练机制}：采用前向传播计算预测值，再通过反向传播算法（基于链式法则）计算损失函数对各层参数的梯度，结合梯度下降法更新权重与偏置，最小化预测损失。
    \item \textbf{特点与局限}：结构简单易实现，理论上可拟合任意非线性函数；但全连接结构参数量大，对图像、序列等结构化数据难以提取空间/时序关联，易出现过拟合。
    \item \textbf{典型应用}：表格数据分类/回归、简单模式识别、作为其他深度模型的输出头模块。
\end{itemize}

\subparagraph{公式解析}
\begin{enumerate}
    \item \textbf{前向传播公式}
    \[
    z^{(l)} = W^{(l)} a^{(l-1)} + b^{(l)},\quad a^{(l)} = \sigma\left(z^{(l)}\right)
    \]
    解析：表示第 $l$ 层的前向计算过程。先通过权重矩阵 $W^{(l)}$ 和偏置 $b^{(l)}$ 对上一层输出 $a^{(l-1)}$ 做线性变换，再经过激活函数 $\sigma(\cdot)$ 得到当前层的激活输出 $a^{(l)}$。其中 $a^{(0)}=x$ 为输入特征，最终输出层结果记为 $\hat{y}=a^{(L)}$，$L$ 为网络总层数。

    \item \textbf{均方误差损失函数}
    \[
    L = \frac{1}{2N}\sum_{i=1}^N \|y_i - \hat{y}_i\|^2
    \]
    解析：回归任务中常用的损失函数，量化模型预测值与真实标签的整体偏差。$N$ 为样本总数，$y_i$ 为真实标签，$\hat{y}_i$ 为模型预测值，对误差取平方后求平均，使误差越大惩罚越重。

    \item \textbf{参数更新规则}
    \[
    W^{(l)} \leftarrow W^{(l)} - \eta \frac{\partial L}{\partial W^{(l)}},\quad b^{(l)} \leftarrow b^{(l)} - \eta \frac{\partial L}{\partial b^{(l)}}
    \]
    解析：基于梯度下降法的参数更新式。$\eta$ 为学习率，控制每次参数更新的步长；$\frac{\partial L}{\partial W^{(l)}}$ 和 $\frac{\partial L}{\partial b^{(l)}}$ 为损失对权重和偏置的梯度，由反向传播计算得到。沿梯度反方向更新参数，可逐步减小损失。
\end{enumerate}

\paragraph{卷积神经网络(CNN)}
\subparagraph{模型概念}
卷积神经网络（Convolutional Neural Network, CNN）是专为处理网格结构数据（如图像）设计的神经网络，核心思想是\textbf{局部感受野}与\textbf{权值共享}，大幅降低参数量并保留空间特征。
\begin{itemize}
    \item \textbf{核心组件}：
    \begin{enumerate}
        \item 卷积层：通过滑动的卷积核在输入特征图上进行局部加权求和，提取边缘、纹理、形状等层级化空间特征；同一卷积核参数共享，显著减少计算量。
        \item 池化层：对特征图进行下采样（如最大池化、平均池化），压缩特征尺寸、保留关键信息，同时提升模型的平移不变性。
        \item 全连接层：在网络末端将二维特征图展平为一维向量，完成最终的分类或回归输出。
    \end{enumerate}
    \item \textbf{特点}：具备强大的空间特征提取能力，参数效率远高于全连接网络，对图像的平移、缩放、旋转具有一定鲁棒性。
    \item \textbf{典型应用}：图像分类、目标检测、图像分割、人脸识别等计算机视觉任务。
\end{itemize}

\subparagraph{公式解析}
\begin{enumerate}
    \item \textbf{二维卷积运算}
    \[
    Y_{i,j} = \sum_{m=0}^{k-1} \sum_{n=0}^{k-1} X_{i+m, j+n} \cdot K_{m,n} + b
    \]
    解析：二维卷积的核心计算式。输入特征图 $X$ 尺寸为 $H\times W$，尺寸为 $k\times k$ 的卷积核 $K$ 在特征图上滑动，对应位置元素相乘求和后加上偏置 $b$，得到输出特征图 $Y$ 的对应元素，实现局部空间特征提取。

    \item \textbf{输出特征图尺寸计算}
    \[
    H_{out} = \frac{H - k + 2P}{S} + 1,\quad W_{out} = \frac{W - k + 2P}{S} + 1
    \]
    解析：用于计算卷积操作后特征图的空间尺寸。$H,W$ 为输入特征图的高和宽，$k$ 为卷积核大小，$P$ 为填充像素数，$S$ 为滑动步长。填充用于控制输出尺寸，步长用于实现下采样。

    \item \textbf{池化运算}
    
    最大池化：
    \[
    Y_{i,j} = \max_{0\leq m,n < k} X_{i\cdot S + m,\ j\cdot S + n}
    \]
    平均池化：
    \[
    Y_{i,j} = \frac{1}{k^2} \sum_{m=0}^{k-1} \sum_{n=0}^{k-1} X_{i\cdot S + m,\ j\cdot S + n}
    \]
    解析：池化用于对特征图进行下采样。最大池化取窗口内最大值，保留最显著的局部特征；平均池化取窗口内平均值，保留整体背景信息。二者均可压缩特征尺寸，并提升模型的平移不变性。
\end{enumerate}

\paragraph{循环神经网络(RNN)}
\subparagraph{模型概念}
循环神经网络（Recurrent Neural Network, RNN）是处理序列数据的专用网络，通过内部的循环状态保留历史信息，使网络具备「记忆」能力，适配文本、语音、时间序列等长度可变的序列输入。
\begin{itemize}
    \item \textbf{核心机制}：每个时间步的神经元不仅接收当前时刻的输入，还接收上一时刻的隐藏状态，从而将历史信息融入当前计算，输出结果由当前输入与历史状态共同决定。
    \item \textbf{固有缺陷}：长序列训练时易出现\textbf{梯度消失}或\textbf{梯度爆炸}问题，难以捕捉长距离依赖关系。
    \item \textbf{主流变体}：长短期记忆网络（LSTM）通过门控机制（输入门、遗忘门、输出门）控制信息的留存与更新；门控循环单元（GRU）简化了门控结构，二者均有效缓解了长序列梯度问题。
    \item \textbf{典型应用}：自然语言处理（语言模型、机器翻译）、语音识别、时间序列预测。
\end{itemize}

\subparagraph{公式解析}
\begin{enumerate}
    \item \textbf{隐藏状态更新公式}
    \[
    h_t = \tanh\left(W_{xh} x_t + W_{hh} h_{t-1} + b_h\right)
    \]
    解析：描述单个时间步的状态更新过程。$x_t$ 为第 $t$ 时刻的输入，$h_{t-1}$ 为上一时刻的隐藏状态，$W_{xh}$ 为输入权重矩阵，$W_{hh}$ 为隐藏状态权重矩阵，$b_h$ 为偏置。经 $\tanh$ 激活后得到当前时刻隐藏状态 $h_t$，实现历史信息与当前输入的融合。

    \item \textbf{输出计算}
    \[
    y_t = W_{hy} h_t + b_y
    \]
    解析：由当前隐藏状态 $h_t$ 经线性变换得到该时间步的输出。$W_{hy}$ 为输出权重矩阵，$b_y$ 为输出偏置；分类任务可在此基础上接入 Softmax 得到类别概率。

    \item \textbf{序列总损失}
    \[
    L = \sum_{t=1}^T L(y_t, \hat{y}_t)
    \]
    解析：序列任务的总损失为所有时间步损失的累加。$T$ 为序列总长度，$L(y_t, \hat{y}_t)$ 为单个时间步的损失函数，训练时通过随时间反向传播（BPTT）算法计算梯度。
\end{enumerate}

\subsubsection{Transformer}
\paragraph{自注意力机制(Self-Attention Mechanism)}
\subparagraph{机制解析}
自注意力机制是Transformer的核心组件，能够让模型在处理序列中每个位置时，动态关注序列中所有相关位置的信息，从而高效捕捉长距离依赖。
\begin{itemize}
    \item \textbf{计算流程}：
    \begin{enumerate}
        \item 将输入序列通过线性变换生成三个向量：查询向量（Query, Q）、键向量（Key, K）、值向量（Value, V）。
        \item 计算Q与所有K的点积，除以维度的平方根进行缩放，再通过Softmax归一化得到注意力权重。
        \item 将注意力权重与对应的V加权求和，得到最终的注意力输出。
    \end{enumerate}
    \item \textbf{核心优势}：并行计算能力远强于RNN，可直接建立序列中任意两个位置的关联，长距离依赖捕捉效率极高。
\end{itemize}

\subparagraph{公式解析}
\begin{enumerate}
    \item \textbf{Q、K、V 线性映射}
    \[
    Q = X W_Q,\quad K = X W_K,\quad V = X W_V
    \]
    解析：将输入序列 $X$ 通过三组可学习的权重矩阵，分别投影为查询向量 $Q$、键向量 $K$、值向量 $V$。$X\in\mathbb{R}^{n\times d_{model}}$，其中 $n$ 为序列长度，$d_{model}$ 为特征维度，$W_Q,W_K,W_V$ 为可训练参数。

    \item \textbf{缩放点积注意力}
    \[
    \text{Attention}(Q,K,V) = \text{softmax}\left( \frac{Q K^T}{\sqrt{d_k}} \right) V
    \]
    解析：自注意力的核心计算公式。先计算 $Q$ 与 $K$ 的点积得到相似度，除以 $\sqrt{d_k}$ 进行缩放，避免点积值过大导致 Softmax 梯度消失；经 Softmax 归一化得到注意力权重，再与 $V$ 加权求和，得到最终的注意力输出。
\end{enumerate}

\paragraph{多头注意力(Multi-Head Attention)}
\subparagraph{机制解析}
多头注意力是对自注意力的扩展，将Q、K、V映射到多个不同的子空间，分别计算注意力后再拼接融合，使模型能从多个维度捕捉特征关联。
\begin{itemize}
    \item \textbf{实现逻辑}：将原始的Q、K、V通过多组独立的线性层拆分为$h$个头（head），每个头独立执行缩放点积注意力计算；最后将所有头的输出拼接，再通过一次线性变换得到最终结果。
    \item \textbf{优势}：不同注意力头可以关注不同类型的依赖关系（如语法关系、语义关联、长/短距离依赖），显著提升模型的特征表达能力。
\end{itemize}

\subparagraph{公式解析}
\begin{enumerate}
    \item \textbf{单头注意力输出}
    \[
    head_i = \text{Attention}\left(Q W_{Q_i},\ K W_{K_i},\ V W_{V_i}\right)
    \]
    解析：将 Q、K、V 分别投影到第 $i$ 个独立的特征子空间，计算该子空间下的注意力结果。每个头对应一组独立的投影矩阵，可捕捉不同类型的特征关联。

    \item \textbf{多头输出拼接与投影}
    \[
    \text{MultiHead}(Q,K,V) = \text{Concat}(head_1, head_2, \dots, head_h) \cdot W_O
    \]
    解析：将 $h$ 个注意力头的输出拼接，再通过输出投影矩阵 $W_O$ 变换为与输入维度一致的结果。多头机制可同时从多个维度捕捉序列依赖，提升模型表达能力。
\end{enumerate}

\paragraph{BERT与GPT模型(BERT \& GPT Models)}
\subparagraph{模型解析}
BERT与GPT是基于Transformer的两大经典预训练语言模型，分别采用编码器和解码器架构，代表了预训练NLP模型的两大主流方向。
\begin{enumerate}
    \item \textbf{BERT（Bidirectional Encoder Representations from Transformers）}：
    \begin{itemize}
        \item 架构：基于Transformer编码器，支持双向注意力，可同时利用上下文信息。
        \item 预训练任务：掩码语言模型（MLM，随机遮盖部分词并预测）、下一句预测（NSP，判断两句话是否连续）。
        \item 特点：双向语义理解能力强，适配分类、实体识别、问答等理解类NLP任务。
    \end{itemize}
    \item \textbf{GPT（Generative Pre-trained Transformer）}：
    \begin{itemize}
        \item 架构：基于Transformer解码器，采用因果注意力（掩码自注意力），仅能利用当前位置及之前的上文信息。
        \item 预训练任务：自回归语言模型，根据上文预测下一个词。
        \item 特点：生成能力突出，适配文本生成、对话、续写等生成类任务，当前主流大语言模型多基于GPT架构扩展。
    \end{itemize}
\end{enumerate}

\subparagraph{公式解析}
\begin{enumerate}
    \item \textbf{GPT 自回归语言模型目标}
    \[
    P(x_{1:T}) = \prod_{t=1}^T P(x_t \mid x_{1:t-1})
    \]
    解析：自回归语言模型的联合概率分解式。将长度为 $T$ 的序列的联合概率，分解为每个位置基于上文的条件概率的乘积。模型仅利用当前位置之前的文本预测下一个词，是生成类语言模型的核心目标。

    \item \textbf{GPT 训练损失}
    \[
    L = -\frac{1}{T} \sum_{t=1}^T \log P(x_t \mid x_{1:t-1})
    \]
    解析：训练损失为负对数似然，通过最大化序列的联合概率，最小化该损失值，使模型生成的文本更符合真实语言分布。

    \item \textbf{BERT 掩码语言模型(MLM)损失}
    \[
    L_{MLM} = -\frac{1}{N_{mask}} \sum_{i \in mask} \log P(x_i \mid x_{\setminus mask})
    \]
    解析：BERT 预训练的核心损失函数。随机掩码序列中部分词元，模型根据未掩码的上下文预测被掩码位置的原始词元。$N_{mask}$ 为掩码词元总数，损失为掩码位置的交叉熵均值，驱动模型学习双向语义表示。
\end{enumerate}

\paragraph{视觉Transformer(Vision Transformer(ViT))}
\subparagraph{模型解析}
视觉Transformer（Vision Transformer, ViT）将Transformer架构迁移到计算机视觉领域，打破了CNN在视觉任务的垄断，在大数据集下表现更优。
\begin{itemize}
    \item \textbf{核心思路}：将二维图像拆分为固定大小的图像块（Patch），将每个图像块展平并映射为一维向量，作为Transformer的输入序列；同时加入可学习的位置编码，保留图像的空间位置信息。
    \item \textbf{结构}：主要由Patch Embedding层、位置编码、多层Transformer编码器、分类头构成；通常加入一个特殊的 \texttt{CLS} token，其最终输出用于图像分类。
    \item \textbf{特点与局限}：全局建模能力强，模型扩展性好，数据量充足时性能优于CNN；但小数据集下表现不如CNN，对局部细节特征的提取弱于卷积。
    \item \textbf{典型应用}：图像分类、目标检测、图像分割等各类视觉任务。
\end{itemize}

\subparagraph{公式解析}
\begin{enumerate}
    \item \textbf{图像块线性嵌入}
    \[
    z_p^i = x_p^i \cdot E,\quad E \in \mathbb{R}^{(P^2 C) \times D}
    \]
    解析：将图像拆分为固定大小的图像块后，每个展平的图像块 $x_p^i$ 通过嵌入矩阵 $E$ 映射为 $D$ 维的特征向量。$P$ 为图像块尺寸，$C$ 为图像通道数，$N=HW/P^2$ 为图像块总数。

    \item \textbf{输入序列构造}
    \[
    z_0 = \left[ x_{class};\ z_p^1;\ z_p^2;\ \dots;\ z_p^N \right] + E_{pos}
    \]
    解析：Transformer 的输入序列由分类 token 与所有图像块嵌入拼接后，加上位置编码构成。$x_{class}$ 为可学习的分类 token，最终用于分类输出；$E_{pos}$ 为位置编码矩阵，用于保留图像的空间位置信息。
\end{enumerate}


\subsubsection{生成模型(Generative Models)}
\paragraph{生成对抗网络(GAN)}
\subparagraph{模型解析}
生成对抗网络（Generative Adversarial Network, GAN）是一种基于博弈思想的生成模型，通过生成器与判别器的动态对抗训练，学习真实数据的分布。
\begin{itemize}
    \item \textbf{双网络结构}：
    \begin{enumerate}
        \item 生成器（Generator）：接收随机噪声作为输入，目标是生成以假乱真的样本，欺骗判别器。
        \item 判别器（Discriminator）：二分类模型，目标是准确区分输入样本是真实数据还是生成器伪造的数据。
    \end{enumerate}
    \item \textbf{训练过程}：二者交替训练，生成器优化生成效果，判别器优化分辨能力，最终达到纳什均衡——生成器生成的样本与真实分布一致，判别器无法区分真假。
    \item \textbf{特点与问题}：生成样本质量高，但训练不稳定，易出现模式崩塌（生成样本多样性不足）、梯度消失等问题。
    \item \textbf{典型应用}：图像生成、图像风格迁移、超分辨率、数据增强。
\end{itemize}

\subparagraph{公式解析}
GAN 的核心是极小极大博弈，目标函数为：
\[
\min_G \max_D V(D,G) = \mathbb{E}_{x\sim p_{data}(x)} \left[ \log D(x) \right] + \mathbb{E}_{z\sim p_z(z)} \left[ \log\left(1 - D(G(z))\right) \right]
\]
解析：该式构成生成器 $G$ 与判别器 $D$ 的二人零和博弈目标。
- $D(x)$ 为判别器判断真实样本 $x$ 为真的概率，取值范围 $[0,1]$；
- $G(z)$ 为生成器将随机噪声 $z$ 映射为生成样本的输出；
- $p_{data}$ 为真实数据分布，$p_z$ 为噪声先验分布，通常为标准正态分布。
训练时固定生成器，最大化判别器的分辨能力；固定判别器，最小化生成样本被判假的概率，二者交替优化直至均衡。

\paragraph{变分自编码器(VAE)}
\subparagraph{模型解析}
变分自编码器（Variational Autoencoder, VAE）是基于变分推断的生成模型，结合了自编码器的结构与概率统计思想，学习数据的潜变量分布。
\begin{itemize}
    \item \textbf{核心结构}：
    \begin{enumerate}
        \item 编码器（推断网络）：将输入数据映射为潜变量空间的均值和方差，拟合后验分布。
        \item 解码器（生成网络）：从潜变量分布中采样，重构出原始数据。
    \end{enumerate}
    \item \textbf{训练目标}：最小化重构误差（保证生成样本与真实样本相似），同时约束潜变量分布趋近于标准正态分布（通过KL散度衡量），即最大化证据下界（ELBO）。
    \item \textbf{特点}：训练稳定，有明确的概率解释，潜空间连续平滑；但生成样本的清晰度通常弱于GAN和扩散模型。
    \item \textbf{典型应用}：数据生成、潜空间插值、异常检测、表示学习。
\end{itemize}

\subparagraph{公式解析}
\begin{enumerate}
    \item \textbf{证据下界（ELBO）}
    \[
    \mathcal{L}(\theta,\phi; x) = \mathbb{E}_{z\sim q_\phi(z|x)} \left[ \log p_\theta(x|z) \right] - D_{KL}\left( q_\phi(z|x) \parallel p(z) \right)
    \]
    解析：VAE 的训练目标为最大化该证据下界，它是对数似然 $\log p_\theta(x)$ 的下界。
    - 第一项为重构项，衡量解码器从潜变量重构原始数据的能力；
    - 第二项为 KL 散度正则项，约束编码器拟合的近似后验 $q_\phi(z|x)$ 趋近于先验分布 $p(z)$，通常为标准正态分布；
    - $\phi$ 为编码器参数，$\theta$ 为解码器参数。

    \item \textbf{高斯假设下的KL散度闭式解}
    \[
    D_{KL}\left(\mathcal{N}(\mu,\sigma^2 I) \parallel \mathcal{N}(0,I)\right) = \frac{1}{2} \sum_{i=1}^d \left( 1 + \log \sigma_i^2 - \mu_i^2 - \sigma_i^2 \right)
    \]
    解析：当近似后验服从对角协方差的高斯分布时，KL 散度可直接通过该式计算，无需采样近似。$d$ 为潜变量维度，$\mu$ 和 $\sigma$ 为编码器输出的潜变量均值和标准差。
\end{enumerate}

\paragraph{扩散模型(Diffusion)}
\subparagraph{模型解析}
扩散模型（Diffusion Model）是一种基于逐步去噪的生成模型，通过正向扩散加噪、反向扩散去噪的过程，学习从随机噪声还原为真实数据的映射。
\begin{itemize}
    \item \textbf{核心过程}：
    \begin{enumerate}
        \item 正向扩散：逐步向真实数据中加入高斯噪声，经过$T$步后数据变为纯随机噪声。
        \item 反向扩散：训练神经网络（通常为U-Net），逐步预测并去除每一步的噪声，从纯噪声中还原出真实数据。
    \end{enumerate}
    \item \textbf{特点}：训练过程稳定，生成样本质量高、多样性好，解决了GAN模式崩塌的问题；但推理速度较慢，需要多步迭代。
    \item \textbf{典型应用}：文生图、图像编辑、视频生成、音频生成，是当前AIGC领域的核心技术之一。
\end{itemize}

\subparagraph{公式解析}
\begin{enumerate}
    \item \textbf{正向扩散加噪公式}
    \[
    x_t = \sqrt{1-\beta_t} \cdot x_{t-1} + \sqrt{\beta_t} \cdot \epsilon_t,\quad \epsilon_t \sim \mathcal{N}(0,I)
    \]
    解析：正向扩散过程逐步向数据中添加高斯噪声。$\beta_t$ 为预设的噪声调度系数，随步数 $t$ 增大单调递增；$\epsilon_t$ 为第 $t$ 步添加的高斯噪声。经过 $T$ 步后，数据趋近于纯随机噪声。

    \item \textbf{任意步加噪结果（重参数化）}
    \[
    x_t = \sqrt{\bar{\alpha}_t} x_0 + \sqrt{1-\bar{\alpha}_t} \epsilon,\quad \alpha_t = 1-\beta_t,\ \bar{\alpha}_t = \prod_{s=1}^t \alpha_s
    \]
    解析：通过重参数化技巧，可直接计算任意第 $t$ 步的加噪结果，无需逐步迭代，大幅提升训练效率。

    \item \textbf{反向去噪训练目标}
    \[
    L = \mathbb{E}_{t, x_0, \epsilon} \left[ \left\| \epsilon - \epsilon_\theta(x_t, t) \right\|^2 \right]
    \]
    解析：扩散模型的训练损失。训练网络 $\epsilon_\theta(x_t, t)$ 预测正向过程中添加的真实噪声 $\epsilon$，损失为预测噪声与真实噪声的均方误差。推理时从纯噪声出发，逐步去噪还原出清晰数据。
\end{enumerate}




%机器学习
\subsection{机器学习(Machine Learning)}
\subsubsection{监督学习(Supervised Learning)}
监督学习使用标记数据进行训练，模型学习输入特征与输出标签之间的映射关系，以对新数据进行预测和分类。
\paragraph{线性回归(Linear Regression)}
\subparagraph{模型解析}
线性回归（Linear Regression）是最基础的监督学习回归算法，通过拟合线性函数描述自变量与因变量之间的关系。
\begin{itemize}
    \item \textbf{模型形式}：假设因变量与自变量满足线性关系 $y = \boldsymbol{w}^T \boldsymbol{x} + b$，其中$\boldsymbol{w}$为权重向量，$b$为偏置项。
    \item \textbf{训练目标}：最小化预测值与真实值的均方误差（MSE），常用求解方法为最小二乘法（闭式解）或梯度下降法。
    \item \textbf{评价指标}：均方误差（MSE）、均方根误差（RMSE）、决定系数$R^2$。
    \item \textbf{特点}：模型简单、可解释性强，计算效率高；但只能拟合线性关系，对异常值敏感。
\end{itemize}

\subparagraph{公式解析}
\begin{enumerate}
    \item \textbf{模型表达式}
    \[
    \hat{y} = \boldsymbol{w}^T \boldsymbol{x} + b
    \]
    解析：线性回归的核心形式，假设因变量与自变量呈线性关系。$\boldsymbol{w}$ 为权重向量，衡量每个特征的重要程度；$b$ 为偏置项，代表基准值；$\boldsymbol{x}$ 为输入特征向量，$\hat{y}$ 为预测值。

    \item \textbf{均方误差损失函数}
    \[
    J(\boldsymbol{w}) = \frac{1}{2n} \| \boldsymbol{y} - \boldsymbol{X}\boldsymbol{w} \|_2^2
    \]
    解析：线性回归的优化目标，最小化预测值与真实值的均方误差。$n$ 为样本数，$\boldsymbol{X}$ 为样本特征矩阵，$\boldsymbol{y}$ 为真实标签向量，系数 $1/2$ 用于求导后消去系数。

    \item \textbf{最小二乘法闭式解}
    \[
    \boldsymbol{w}^* = (\boldsymbol{X}^T \boldsymbol{X})^{-1} \boldsymbol{X}^T \boldsymbol{y}
    \]
    解析：对损失函数求导并令导数为零，直接得到的最优参数解析解。当特征维度不高时可直接计算，无需迭代优化。

    \item \textbf{梯度下降更新规则}
    \[
    \boldsymbol{w} \leftarrow \boldsymbol{w} - \eta \cdot \frac{1}{n} \boldsymbol{X}^T (\boldsymbol{X}\boldsymbol{w} - \boldsymbol{y})
    \]
    解析：梯度下降法的参数更新式，通过迭代逐步逼近最优解。$\eta$ 为学习率，控制每次更新的步长，适合样本量较大的场景。
\end{enumerate}

\paragraph{逻辑回归(Logistic Regression)}
\subparagraph{模型解析}
逻辑回归（Logistic Regression）是经典的二分类算法，在线性回归基础上引入Sigmoid激活函数，将线性输出映射为0~1之间的概率值。
\begin{itemize}
    \item \textbf{核心逻辑}：通过 $h_\theta(\boldsymbol{x}) = \sigma(\boldsymbol{w}^T \boldsymbol{x} + b)$ 计算样本属于正类的概率，其中$\sigma(z)=\frac{1}{1+e^{-z}}$为Sigmoid函数；通常设定阈值（如0.5）进行类别判定。
    \item \textbf{训练目标}：最大化对数似然函数，等价于最小化交叉熵损失函数，通常用梯度下降法求解。
    \item \textbf{特点}：模型简单、训练快、可解释性强（可通过权重分析特征重要性）；但只能处理线性可分问题，对非线性数据拟合能力有限。
    \item \textbf{扩展}：可通过Softmax函数扩展为多分类逻辑回归。
\end{itemize}

\subparagraph{公式解析}
\begin{enumerate}
    \item \textbf{Sigmoid函数}
    \[
    \sigma(z) = \frac{1}{1 + e^{-z}},\quad z = \boldsymbol{w}^T \boldsymbol{x} + b
    \]
    解析：将线性输出 $z$ 映射到 $(0,1)$ 区间，作为样本属于正类的概率。函数单调递增，输出具有概率意义，是逻辑回归的核心映射函数。

    \item \textbf{交叉熵损失函数}
    \[
    J(\boldsymbol{w}) = -\frac{1}{n} \sum_{i=1}^n \left[ y_i \log \hat{y}_i + (1-y_i) \log (1-\hat{y}_i) \right]
    \]
    解析：逻辑回归的优化目标，即负对数似然损失。$y_i$ 为真实标签（0或1），$\hat{y}_i=\sigma(z_i)$ 为预测正类概率。预测与真实标签越接近，损失值越小。

    \item \textbf{梯度更新公式}
    \[
    \boldsymbol{w} \leftarrow \boldsymbol{w} - \eta \cdot \frac{1}{n} \sum_{i=1}^n (\hat{y}_i - y_i) \boldsymbol{x}_i
    \]
    解析：交叉熵损失对权重的梯度形式简洁，通过梯度下降迭代更新参数，最小化分类损失。
\end{enumerate}

\paragraph{决策树(Decision Tree)}
\subparagraph{模型解析}
决策树（Decision Tree）是一种基于树结构的监督学习算法，通过递归地对特征进行分裂，实现样本的分类或回归预测。
\begin{itemize}[leftmargin=*]
    \item \textbf{树结构}：根节点包含全部样本，内部节点对应一个特征分裂条件，叶节点对应最终的预测结果。
    \item \textbf{分裂准则}：分类任务常用信息增益（ID3）、信息增益比（C4.5）、基尼系数（CART）；回归任务常用平方误差最小化。
    \item \textbf{训练过程}：递归选择最优分裂特征，将数据集划分为子集，直到满足停止条件（如节点样本数过少、信息增益低于阈值、达到最大深度）。
    \item \textbf{特点}：可解释性强（可生成可视化规则），无需特征归一化，能处理离散与连续特征；但易过拟合，对数据波动敏感。
    \item \textbf{优化方向}：剪枝（预剪枝、后剪枝）降低过拟合风险。
\end{itemize}

\subparagraph{公式解析}
\begin{enumerate}
    \item \textbf{信息熵}
    \[
    H(D) = -\sum_{k=1}^K p_k \log_2 p_k
    \]
    解析：衡量数据集 $D$ 的纯度，$K$ 为类别总数，$p_k$ 为第 $k$ 类样本的占比。熵值越大，数据集混乱程度越高；熵值为 0 时数据集类别完全统一。

    \item \textbf{信息增益}
    \[
    g(D,A) = H(D) - H(D|A)
    \]
    解析：按特征 $A$ 分裂数据集后，信息熵的减少量，即纯度的提升量。$H(D|A)$ 为分裂后的条件熵。信息增益越大，说明该特征的分类效果越好，是 ID3 算法的分裂准则。

    \item \textbf{基尼系数}
    \[
    Gini(D) = 1 - \sum_{k=1}^K p_k^2
    \]
    解析：表示从数据集中随机抽取两个样本，类别不一致的概率。基尼系数越小，数据集纯度越高，是 CART 分类树的分裂准则，计算效率高于信息熵。
\end{enumerate}

\paragraph{朴素贝叶斯(Naive Bayes)}
\subparagraph{模型解析}
朴素贝叶斯（Naive Bayes）是基于贝叶斯定理与\textbf{特征条件独立假设}的分类算法，属于生成式模型。
\begin{itemize}
    \item \textbf{核心公式}：基于贝叶斯定理 $P(y|\boldsymbol{x}) = \frac{P(\boldsymbol{x}|y)P(y)}{P(\boldsymbol{x})}$，其中$P(y)$为先验概率，$P(\boldsymbol{x}|y)$为类条件概率。「朴素」即假设所有特征之间相互独立，因此$P(\boldsymbol{x}|y)=\prod_{i=1}^n P(x_i|y)$。
    \item \textbf{常见变体}：高斯朴素贝叶斯（处理连续特征，假设特征服从正态分布）、多项式朴素贝叶斯（处理离散计数特征，如文本词频）、伯努利朴素贝叶斯（处理二值特征）。
    \item \textbf{特点}：算法简单、训练速度极快，对小数据友好，适合高维数据；但特征独立假设在实际场景中往往不成立，可能影响分类精度。
    \item \textbf{典型应用}：文本分类、垃圾邮件识别、情感分析。
\end{itemize}

\subparagraph{公式解析}
\begin{enumerate}
    \item \textbf{贝叶斯定理}
    \[
    P(y_k | \boldsymbol{x}) = \frac{P(\boldsymbol{x} | y_k) P(y_k)}{P(\boldsymbol{x})}
    \]
    解析：朴素贝叶斯的核心公式，由先验概率 $P(y_k)$ 和类条件概率 $P(\boldsymbol{x}|y_k)$ 计算后验概率 $P(y_k|\boldsymbol{x})$。$P(\boldsymbol{x})$ 为证据因子，对所有类别取值相同，分类时可忽略。

    \item \textbf{特征条件独立假设}
    \[
    P(\boldsymbol{x} | y_k) = \prod_{i=1}^d P(x_i | y_k)
    \]
    解析：“朴素”的核心假设，类别确定时各特征相互独立，将高维的类条件概率分解为各特征概率的乘积，大幅简化计算。$d$ 为特征维度。

    \item \textbf{分类决策规则}
    \[
    \hat{y} = \arg\max_{y_k} P(y_k) \prod_{i=1}^d P(x_i | y_k)
    \]
    解析：选取后验概率最大的类别作为预测结果，无需计算分母的证据项，直接比较分子部分即可。
\end{enumerate}

\subsubsection{高级监督学习(Advanced Supervised Learning)}
\paragraph{支持向量机(SVM)}
\subparagraph{模型解析}
支持向量机（Support Vector Machine, SVM）是一种经典的二分类算法，核心目标是找到一个最大间隔超平面，将不同类别的样本分隔开，使模型泛化能力最优。
\begin{itemize}
    \item \textbf{核心概念}：
    \begin{enumerate}
        \item 间隔：超平面到两侧最近样本的距离之和，最大化间隔可提升模型鲁棒性。
        \item 支持向量：距离超平面最近的样本点，决定了超平面的位置，是模型的关键参数。
    \end{enumerate}
    \item \textbf{核技巧}：对于线性不可分数据，通过核函数将低维特征映射到高维空间，使其在高维空间线性可分；常用核函数包括线性核、多项式核、高斯核（RBF）。
    \item \textbf{软间隔}：允许部分样本不满足间隔约束，通过惩罚参数C平衡间隔大小与错分样本数，缓解过拟合。
    \item \textbf{特点}：小样本下泛化能力强，适合高维数据，核技巧可有效处理非线性问题；但大数据集训练慢，对参数和核函数选择敏感。
\end{itemize}

\subparagraph{公式解析}
\begin{enumerate}
    \item \textbf{硬间隔SVM原始优化问题}
    \[
    \begin{aligned}
    \min_{\boldsymbol{w}, b} \quad & \frac{1}{2} \| \boldsymbol{w} \|^2 \\
    \text{s.t.} \quad & y_i (\boldsymbol{w}^T \boldsymbol{x}_i + b) \geq 1,\quad i=1,2,\dots,n
    \end{aligned}
    \]
    解析：硬间隔支持向量机的优化目标。最小化 $\|\boldsymbol{w}\|^2$ 等价于最大化分类间隔；约束条件保证所有样本都被正确分类且满足间隔要求。$y_i\in\{-1,+1\}$ 为样本标签。

    \item \textbf{软间隔SVM}
    \[
    \begin{aligned}
    \min_{\boldsymbol{w}, b, \xi} \quad & \frac{1}{2} \| \boldsymbol{w} \|^2 + C \sum_{i=1}^n \xi_i \\
    \text{s.t.} \quad & y_i (\boldsymbol{w}^T \boldsymbol{x}_i + b) \geq 1 - \xi_i \\
    & \xi_i \geq 0
    \end{aligned}
    \]
    解析：引入松弛变量 $\xi_i$，允许部分样本不满足间隔约束，适配线性不可分的数据。惩罚参数 $C$ 用于平衡间隔大小与错分样本的数量，$C$ 越大对错分的惩罚越重。
\end{enumerate}

\paragraph{随机森林(Random Forest)}
\subparagraph{模型解析}
随机森林（Random Forest）是基于Bagging集成思想的算法，由多棵决策树组成，通过「数据随机+特征随机」降低单棵树的过拟合，提升整体稳定性与精度。
\begin{itemize}
    \item \textbf{集成逻辑}：
    \begin{enumerate}
        \item 数据随机：对训练集进行有放回抽样（Bootstrap采样），生成多个不同的子数据集，分别训练一棵决策树。
        \item 特征随机：每棵树分裂节点时，仅从随机选取的部分特征中选择最优分裂特征。
    \end{enumerate}
    \item \textbf{预测方式}：分类任务采用多数投票法，回归任务采用平均值法。
    \item \textbf{特点}：抗过拟合能力强，对异常值和噪声鲁棒，并行训练效率高，无需复杂调参；但可解释性弱于单棵决策树，小数据集上提升有限。
\end{itemize}

\subparagraph{公式解析}
\begin{enumerate}
    \item \textbf{分类任务：多数投票法}
    \[
    H(\boldsymbol{x}) = \arg\max_{y \in \mathcal{Y}} \sum_{t=1}^T \mathbb{I} \left( h_t(\boldsymbol{x}) = y \right)
    \]
    解析：随机森林分类的输出规则。$T$ 为决策树的数量，$h_t(\boldsymbol{x})$ 为第 $t$ 棵树的预测结果，$\mathbb{I}(\cdot)$ 为指示函数，条件成立时取 1。最终输出得票最多的类别。

    \item \textbf{回归任务：均值法}
    \[
    H(\boldsymbol{x}) = \frac{1}{T} \sum_{t=1}^T h_t(\boldsymbol{x})
    \]
    解析：随机森林回归的输出规则，取所有决策树预测结果的平均值作为最终输出，通过平均降低单棵树的方差，提升模型稳定性。
\end{enumerate}

\paragraph{梯度提升(Gradient Boosting)}
\subparagraph{模型解析}
梯度提升（Gradient Boosting）是基于Boosting集成思想的算法，通过串行训练多个弱学习器（通常为决策树），每一轮新的学习器拟合上一轮模型的残差（负梯度），逐步降低预测误差。
\begin{itemize}
    \item \textbf{核心思想}：将多个弱模型加权组合，形成强模型；每轮训练沿着损失函数的负梯度方向优化，不断逼近真实值。
    \item \textbf{经典实现}：
    \begin{enumerate}
        \item GBDT：基础的梯度提升决策树，以CART树为基学习器。
        \item XGBoost：对GBDT的工程优化，加入正则化、二阶泰勒展开、并行计算等，精度与效率更高。
        \item LightGBM：采用直方图算法、按叶子生长策略，训练速度更快，内存占用更低。
    \end{enumerate}
    \item \textbf{特点}：预测精度高，适合结构化数据竞赛；但串行训练无法完全并行，对异常值较敏感，调参相对复杂。
\end{itemize}

\subparagraph{公式解析}
\begin{enumerate}
    \item \textbf{加法模型形式}
    \[
    F_M(\boldsymbol{x}) = \sum_{m=1}^M \gamma_m h_m(\boldsymbol{x})
    \]
    解析：梯度提升模型的整体形式，由 $M$ 个基学习器加权相加得到最终强模型。$h_m(\boldsymbol{x})$ 为第 $m$ 个基学习器，通常为决策树；$\gamma_m$ 为对应学习率。

    \item \textbf{前向分步优化}
    \[
    F_m(\boldsymbol{x}) = F_{m-1}(\boldsymbol{x}) + \gamma_m h_m(\boldsymbol{x})
    \]
    解析：梯度提升的串行训练方式，每一步在现有模型基础上添加一个新的基学习器，逐步降低整体损失。

    \item \textbf{负梯度拟合}
    \[
    r_{mi} = -\left[ \frac{\partial L(y_i, F(\boldsymbol{x}_i))}{\partial F(\boldsymbol{x}_i)} \right]_{F=F_{m-1}}
    \]
    解析：第 $m$ 轮训练的拟合目标，即损失函数对当前模型输出的负梯度，可视为残差的推广形式。新的基学习器拟合该负梯度，沿损失下降方向优化模型。
\end{enumerate}

\subsubsection{无监督学习(Unsupervised Learning)}
无监督学习处理未标记数据，旨在发现数据中的潜在模式、结构或关系，而不需要明确的目标变量或标签。
\paragraph{K均值聚类(K-Means)}
\subparagraph{算法解析}
K-Means是最经典的基于划分的无监督聚类算法，目标是将$n$个样本划分为$K$个簇，使簇内样本相似度高、簇间相似度低。
\begin{itemize}
    \item \textbf{算法流程}：
    \begin{enumerate}
        \item 随机选取$K$个样本作为初始簇中心。
        \item 分配阶段：计算每个样本到所有簇中心的距离，将样本归入距离最近的簇。
        \item 更新阶段：重新计算每个簇的均值，作为新的簇中心。
        \item 重复2、3步，直到簇中心不再变化或达到最大迭代次数。
    \end{enumerate}
    \item \textbf{距离度量}：常用欧氏距离，也可使用曼哈顿距离、余弦相似度等。
    \item \textbf{K值选择}：常用肘部法则（根据簇内误差平方和随K变化的拐点）、轮廓系数法确定最优K。
    \item \textbf{特点}：原理简单、收敛快、适合大数据集；但对初始簇中心敏感，易陷入局部最优，对异常值敏感，需预先指定K值。
\end{itemize}

\subparagraph{公式解析}
\begin{enumerate}
    \item \textbf{优化目标（簇内误差平方和SSE）}
    \[
    J = \sum_{k=1}^K \sum_{\boldsymbol{x} \in C_k} \| \boldsymbol{x} - \boldsymbol{\mu}_k \|_2^2
    \]
    解析：K-Means 的优化目标，最小化所有样本到其所属簇质心的距离平方和。$C_k$ 为第 $k$ 个簇，$\boldsymbol{\mu}_k$ 为第 $k$ 个簇的质心，$K$ 为预设的簇数。

    \item \textbf{簇质心计算}
    \[
    \boldsymbol{\mu}_k = \frac{1}{|C_k|} \sum_{\boldsymbol{x} \in C_k} \boldsymbol{x}
    \]
    解析：每个簇的质心为该簇所有样本的均值，是簇内所有点的中心位置。
\end{enumerate}

\paragraph{层次聚类(Hierarchical Clustering)}
\subparagraph{算法解析}
层次聚类（Hierarchical Clustering）是通过层级方式构建聚类结构的算法，无需预先指定簇数，最终生成树状的聚类谱系图。
\begin{itemize}
    \item \textbf{两类思路}：
    \begin{enumerate}
        \item 凝聚式层次聚类（自底向上）：初始每个样本为一个簇，逐步合并距离最近的两个簇，直到所有样本归为一个簇。
        \item 分裂式层次聚类（自顶向下）：初始所有样本为一个簇，逐步分裂为更小的簇，直到每个样本为一个簇。
    \end{enumerate}
    \item \textbf{簇间距离度量}：最短距离法（单链接）、最长距离法（全链接）、平均距离法、沃德法（最小化簇内方差）。
    \item \textbf{特点}：无需指定簇数，聚类结果层次清晰，可解释性强；但计算复杂度高，不适合大数据集，一旦合并/分裂无法修正。
\end{itemize}

\subparagraph{公式解析}
常用簇间距离度量公式（凝聚式聚类）：
\begin{enumerate}
    \item 单链接（最短距离）
    \[
    d_{min}(C_i, C_j) = \min_{\boldsymbol{x} \in C_i, \boldsymbol{y} \in C_j} \|\boldsymbol{x} - \boldsymbol{y}\|
    \]
    解析：两个簇之间的距离定义为两簇中最近样本的距离，只要两个簇有样本接近就倾向于合并，容易形成链式结构。

    \item 全链接（最长距离）
    \[
    d_{max}(C_i, C_j) = \max_{\boldsymbol{x} \in C_i, \boldsymbol{y} \in C_j} \|\boldsymbol{x} - \boldsymbol{y}\|
    \]
    解析：两个簇之间的距离定义为两簇中最远样本的距离，只有当两个簇所有样本都接近时才合并，对异常值更敏感。

    \item 平均链接
    \[
    d_{avg}(C_i, C_j) = \frac{1}{|C_i||C_j|} \sum_{\boldsymbol{x} \in C_i} \sum_{\boldsymbol{y} \in C_j} \|\boldsymbol{x} - \boldsymbol{y}\|
    \]
    解析：两个簇之间的距离为所有样本对距离的平均值，综合考虑了簇内所有样本的信息，结果更稳定。
\end{enumerate}

\paragraph{主成分分析(PCA)}
\subparagraph{算法解析}
主成分分析（Principal Component Analysis, PCA）是最常用的线性降维算法，通过正交变换将高维相关特征转换为低维互不相关的主成分，同时尽可能保留原始数据的方差信息。
\begin{itemize}
    \item \textbf{核心原理}：寻找数据中方差最大的方向作为第一主成分，第二主成分与第一主成分正交且方差次大，以此类推，最终选取前$k$个主成分实现降维。
    \item \textbf{计算步骤}：数据标准化→计算协方差矩阵→对协方差矩阵做特征值分解→选取前$k$大特征值对应的特征向量→将原始数据投影到特征向量构成的新空间。
    \item \textbf{特点}：计算简单、可解释性较强，能去除数据冗余与噪声；但降维后特征物理含义模糊，且仅能捕捉线性关系，对非线性数据效果差。
    \item \textbf{典型应用}：数据可视化、特征压缩、去噪、加速模型训练。
\end{itemize}

\subparagraph{公式解析}
\begin{enumerate}
    \item \textbf{协方差矩阵}
    \[
    \boldsymbol{\Sigma} = \frac{1}{n-1} \boldsymbol{X}^T \boldsymbol{X}
    \]
    解析：去中心化后的样本矩阵 $\boldsymbol{X}$ 对应的协方差矩阵，描述不同特征之间的线性相关程度。$n$ 为样本数，$\boldsymbol{\Sigma}$ 为 $d\times d$ 的对称矩阵，$d$ 为特征维度。

    \item \textbf{特征值分解}
    \[
    \boldsymbol{\Sigma} \boldsymbol{v}_i = \lambda_i \boldsymbol{v}_i
    \]
    解析：对协方差矩阵做特征值分解，$\lambda_i$ 为特征值，对应主成分的方差大小；$\boldsymbol{v}_i$ 为对应的单位正交特征向量，即主成分的方向。

    \item \textbf{降维投影}
    \[
    \boldsymbol{Y} = \boldsymbol{X} \boldsymbol{W}
    \]
    解析：将特征值从大到小排序，选取前 $k$ 个特征向量构成投影矩阵 $\boldsymbol{W}$，原始数据投影后得到降维结果 $\boldsymbol{Y}$，保留了数据中方差最大的 $k$ 个方向的信息。
\end{enumerate}

\paragraph{t-分布随机邻域嵌入(t-SNE)}
\subparagraph{算法解析}
t-SNE（t-Distributed Stochastic Neighbor Embedding）是非线性降维算法，核心是将高维空间中样本的相似性概率分布，映射到低维空间并保持分布一致，主要用于高维数据可视化。
\begin{itemize}
    \item \textbf{核心思路}：
    \begin{enumerate}
        \item 在高维空间中，用高斯分布描述样本间的相似性概率。
        \item 在低维空间中，用t分布描述样本间的相似性概率（缓解高维数据拥挤问题）。
        \item 通过最小化两个分布的KL散度，优化低维空间的样本坐标。
    \end{enumerate}
    \item \textbf{特点}：能很好地保留数据的局部结构，可视化效果优秀，适合展现数据的簇状分布；但计算复杂度高、速度慢，结果具有随机性，无法直接用于特征工程（仅适合可视化）。
\end{itemize}

\subparagraph{公式解析}
\begin{enumerate}
    \item \textbf{高维空间相似性概率}
    \[
    p_{j|i} = \frac{\exp\left( -\|x_i - x_j\|^2 / 2\sigma_i^2 \right)}{\sum_{k \neq i} \exp\left( -\|x_i - x_k\|^2 / 2\sigma_i^2 \right)}
    \]
    解析：高维空间中，以 $x_i$ 为中心，用高斯分布定义其他样本是其邻域点的条件概率。$\sigma_i$ 为带宽参数，控制邻域范围。对称化后得到联合概率 $p_{ij}$。

    \item \textbf{低维空间相似性概率}
    \[
    q_{ij} = \frac{\left( 1 + \|y_i - y_j\|^2 \right)^{-1}}{\sum_{k \neq l} \left( 1 + \|y_k - y_l\|^2 \right)^{-1}}
    \]
    解析：低维空间中用自由度为 1 的 t 分布定义样本间的相似性概率，长尾特性可缓解高维数据在低维的拥挤问题，保留局部邻域结构。

    \item \textbf{优化目标：KL散度最小化}
    \[
    C = KL(P \parallel Q) = \sum_i \sum_j p_{ij} \log \frac{p_{ij}}{q_{ij}}
    \]
    解析：t-SNE 的优化目标为最小化高维分布 $P$ 与低维分布 $Q$ 的 KL 散度，使低维嵌入尽可能保留高维空间的局部邻域关系，通过梯度下降优化低维坐标。
\end{enumerate}


\end{document}